\section{提案手法：危険度連動eHMI}
\label{sec:method}

\subsection{システム概要}

各AGVに対し，毎フレーム以下の処理を行う（\Cref{fig:system}参照）．

\begin{enumerate}[leftmargin=2em]
  \item 参加者の視線方向から\textbf{動的基準線}（太線ゾーン）を更新
  \item AGVの残存計画経路と基準線の交差から\textbf{TTC}を算出
  \item TTCと近接距離から危険度スコア $R \in [0,1]$ を計算
  \item $R$をHSV色空間の色相・彩度・明度および4段階不透明度に変換
  \item 走行中は残存経路，停止中は停止線として描画
\end{enumerate}

\begin{figure}[htbp]
  \centering
  \fbox{\parbox{0.92\linewidth}{%
    \small
    \textbf{[図1: システム構成（ここに図を挿入）]}\\[0.5em]
    参加者 $\rightarrow$ DynamicCrossingLineTracker（視線基準線）\\
    AGV $\rightarrow$ VehicleRiskCalculator（$R$）\\
    $\rightarrow$ RiskToVisualMapper $\rightarrow$ PathRenderer
  }}
  \caption{提案eHMIの処理パイプライン}
  \label{fig:system}
\end{figure}

\subsection{動的基準線（視線方向の太線ゾーン）}

従来のNavMesh最短経路に基づく基準線ではなく，
参加者が\textbf{現在見ている方向}に注目領域を定義する．

HMDのforwardをXZ平面に投影した単位ベクトル $\mathbf{f}$ に対し，
視線起点 $\mathbf{o}$ からレイキャストを行い，障害物（棚・壁）との
交点 $\mathbf{h}$ を求める．
床面上の基準軸は
\begin{equation}
  \mathbf{a}_\text{start} = \mathbf{p}, \quad
  \mathbf{a}_\text{end} = \mathrm{flatten}(\mathbf{h})
\end{equation}
とする．$\mathbf{p}$ は参加者位置，半幅 $w=\SI{0.5}{m}$（全幅\SI{1.0}{m}）の
帯状ゾーンとして，AGV経路との交差判定に用いる．

この設計により，作業員が立ち止まり別方向を見た場合も，
\textbf{視線に追従した危険評価}が可能となる．

\subsection{危険度スコア}

\subsubsection{表示可否}

まず描画対象か判定する．
\begin{equation}
  \text{visible} = (T \le T_\text{max}) \vee (d \le D_\text{max})
\end{equation}
ここで $T$ はTTC，$d$ はAGVと参加者の3次元距離，
$T_\text{max}=\SI{4.0}{s}$，$D_\text{max}=\SI{13.6}{m}$ である．
表示距離上限は
\begin{equation}
  D_\text{max} = (v_\text{AGV,max} + v_\text{ped}) \times T_\text{max}
  = (2.0 + 1.4) \times 4 = \SI{13.6}{m}
\end{equation}
として，AGV最大速度 $v_\text{AGV,max}=\SI{2.0}{m/s}$，
歩行速度 $v_\text{ped}=\SI{1.4}{m/s}$ から導出する．

\subsubsection{TTC}

AGVの残存計画経路 $\{\mathbf{w}_0,\ldots,\mathbf{w}_{m-1}\}$ 上を
現在速度 $v$ で走行したとき，基準線ゾーンとの最初の交差までの
累積距離 $s$ から
\begin{equation}
  T = \frac{s}{\max(v,\, v_\text{min})}, \quad v_\text{min}=\SI{0.1}{m/s}
\end{equation}
と定義する．交差が $T_\text{max}$ 以内に見つからない場合 $T=\infty$ とする．

\subsubsection{TTC・近接由来スコアと統合}

\begin{align}
  R_\text{ttc} &=
  \begin{cases}
    0 & T = \infty \\
    \left(1 - \dfrac{T}{T_\text{max}}\right)^{\gamma} & 0 \le T \le T_\text{max}
  \end{cases} \\[6pt]
  R_\text{prox} &=
  \begin{cases}
    0 & d \ge D_\text{max} \\
    \left(1 - \dfrac{d}{D_\text{max}}\right)^{\gamma} & d < D_\text{max}
  \end{cases} \\[6pt]
  R &= \max(R_\text{ttc},\, R_\text{prox},\, R_\text{floor})
\end{align}
非線形係数 $\gamma=0.6$，下限 $R_\text{floor}=0.08$ とする．

\subsection{視覚マッピング}

危険度 $R$ をHSV色空間へ写像する（Proposed条件のみ）．
\begin{align}
  H(R) &= (1-R) \times 180° \quad \text{（青緑$\rightarrow$赤）} \\
  S(R) &= 0.4 + 0.6R \\
  V(R) &= 0.5 + 0.3R
\end{align}
不透明度は認知しやすさのため4段階に量子化する．
\begin{equation}
  \alpha(R) \in \{0.35,\, 0.55,\, 0.80,\, 1.00\}
\end{equation}

\subsection{3条件の対照}

\Cref{tab:conditions}に実験条件を示す．
Baselineは「情報あり・危険度なし」，No-ARは「情報なし」，
Proposedは「情報あり・危険度連動」として，被験者内の3水準を構成する．

\begin{table}[htbp]
  \centering
  \caption{実験条件の対照}
  \label{tab:conditions}
  \begin{tabular}{lccc}
    \toprule
    項目 & Baseline & No-AR & Proposed \\
    \midrule
    経路表示 & あり & なし & あり \\
    危険度計算 & なし（$R=1$固定） & --- & TTC + 近接 \\
    色 & 固定青 & --- & HSV連動 \\
    不透明度 & 1.0固定 & --- & 4段階 \\
    表示距離 & \SI{13.6}{m}以内 & --- & \SI{13.6}{m}以内 \\
    \bottomrule
  \end{tabular}
\end{table}

\subsection{経路描画}

走行中は残存計画経路をテーパー付きライン（根元\SI{0.65}{m}，先端\SI{0.12}{m}）で描画する．
停止中（荷物ピックアップ／ドロップ）は停止線（全幅\SI{1.0}{m}）を表示する．
危険度の高いAGVほど手前に描画されるよう，$\mathrm{sortingOrder}=\mathrm{round}(R \times 100)$ とする．

\begin{table}[htbp]
  \centering
  \caption{主要パラメータ一覧}
  \label{tab:params}
  \begin{tabular}{lll}
    \toprule
    記号 & 値 & 意味 \\
    \midrule
    $T_\text{max}$ & \SI{4.0}{s} & TTC表示上限 \\
    $D_\text{max}$ & \SI{13.6}{m} & 近接表示上限 \\
    $\gamma$ & 0.6 & 危険度カーブの急峻さ \\
    $R_\text{floor}$ & 0.08 & 表示中の最低危険度 \\
    $w$ & \SI{0.5}{m} & 基準線半幅 \\
    \bottomrule
  \end{tabular}
\end{table}
